% paper/sections/appendix_a_nondim_details.tex
%
% 付録A: 無次元化の項別導出（§2c から移動）
%   本編 §~\ref{sec:nondim} の式~\eqref{eq:NS_dimless_cls} を導く際の
%   各項の変換過程を詳述する．

\section{無次元 Navier--Stokes 方程式の項別導出}
\label{app:nondim_details}

本付録では，§~\ref{sec:nondim} で提示した無次元 NS 方程式~\eqref{eq:NS_dimless_cls}
の導出過程を，各項ごとに詳述する．

% ─────────────────────────────────────────────────────────────────────────────
\subsection{各項の無次元化}
\label{sec:nondim_items}
% ─────────────────────────────────────────────────────────────────────────────

§~\ref{sec:nondim} で導入した無次元化変数を，
次元系方程式~\eqref{eq:NS_dim} の各項に代入し，変換する．
各項に慣性力スケール $\rho_l U^2 / L$ が現れることを確認する．

\begin{itemize}
    \item \textbf{非定常項（時間微分項）}:
    \[
      \rho^* \pd{\bu^*}{t^*}
      = (\rho_l\rho)\,\pd{(U\bu)}{\frac{L}{U}t}
      = \left[\frac{\rho_l U^2}{L}\right]\,\rho\,\pd{\bu}{t}
    \]

    \item \textbf{移流項}:
    \[
      \rho^* (\bu^* \cdot \bnabla^*) \bu^*
      = (\rho_l\rho) \left( (U\bu) \cdot \left(\frac{1}{L}\bnabla\right) \right) (U\bu)
      = \left[\frac{\rho_l U^2}{L}\right] \rho(\bu\cdot\bnabla)\bu
    \]

    \item \textbf{圧力勾配項}:
    \[
      \bnabla^* p^* = \left(\frac{1}{L}\bnabla\right) (\rho_l U^2 p)
      = \left[\frac{\rho_l U^2}{L}\right] \bnabla p
    \]

    \item \textbf{粘性項}:
    \begin{align*}
     \bnabla^*\cdot\bigl[\mu^*(\bnabla^*\bu^* + (\bnabla^*\bu^*)^T)\bigr]
     &= \frac{1}{L}\bnabla\cdot\left[
         (\mu_l\mu)\left(\frac{U}{L}\bnabla\bu + \frac{U}{L}(\bnabla\bu)^T\right)
        \right] \\
     &= \left[\frac{\mu_l U}{L^2}\right]\bnabla\cdot\bigl[\mu(\bnabla\bu+(\bnabla\bu)^T)\bigr] \\
     &= \left[\frac{\rho_l U^2}{L}\right]
        \underbrace{\left(\frac{\mu_l}{\rho_l U L}\right)}_{=1/Re}
        \bnabla\cdot\bigl[\mu(\bnabla\bu+(\bnabla\bu)^T)\bigr]
    \end{align*}
    ここで，レイノルズ数 $Re \equiv \rho_l U L / \mu_l$ が現れる．

    \item \textbf{重力項}:（座標規約 \eqref{eq:gravity_convention} より $\bm{g}^* = -g\hat{\bm{z}}$）
    \[
      \rho^* \bm{g}^*
      = (\rho_l\rho) (-g\hat{\bm{z}})
      = -\left[\frac{\rho_l U^2}{L}\right]
        \underbrace{\left(\frac{gL}{U^2}\right)}_{=1/Fr^2}
        \rho \hat{\bm z}
    \]
    ここで，フルード数 $Fr \equiv U/\sqrt{gL}$ が現れる．

    \item \textbf{表面張力項}:（無次元数 $We$ の導出には CSF 体積力形式 $\bm{f}_\sigma^* = \sigma\kappa^*\bnabla^*\psi$ を用いるが，得られる $1/We$ は Young--Laplace ジャンプ条件~\eqref{eq:young_laplace_jump} にも共通する）
    \begin{align*}
      \sigma\kappa^*\bnabla^*\psi
      &= \sigma \left(\frac{1}{L}\kappa\right) \left(\frac{1}{L}\bnabla\psi\right)
       = \frac{\sigma}{L^2}\kappa\bnabla\psi \\
      &= \left[\frac{\rho_l U^2}{L}\right]
         \underbrace{\left(\frac{\sigma}{\rho_l U^2 L}\right)}_{=1/We}
         \kappa\bnabla\psi
    \end{align*}
    ここで，ウェーバー数 $We \equiv \rho_l U^2 L / \sigma$ が現れる．
\end{itemize}

上記の各項を元の方程式~\eqref{eq:NS_dim} に代入すると，
共通因子 $[\rho_l U^2/L]$ が全ての項に現れ，
両辺をこの因子で割ることで式~\eqref{eq:NS_dimless_cls} が得られる．

% ─────────────────────────────────────────────────────────────────────────────
\subsection{2 次元曲率式の展開導出}
\label{app:curvature_2d_derivation}
% ─────────────────────────────────────────────────────────────────────────────

§~\ref{sec:curvature} の曲率式~\eqref{eq:curvature_2d} の導出過程を示す．
式~\eqref{eq:curvature_div} を2次元 $(x,y)$ 系で展開する．
$\bnabla\phi = (\phi_x, \phi_y)$，
$|\bnabla\phi| = \sqrt{\phi_x^2 + \phi_y^2}$ として：
\[
  \kappa_{lg} = \left( \pd{n_x}{x} + \pd{n_y}{y} \right)
  = \left(
      \pd{}{x}\frac{\phi_x}{\sqrt{\phi_x^2+\phi_y^2}}
      + \pd{}{y}\frac{\phi_y}{\sqrt{\phi_x^2+\phi_y^2}}
    \right)
\]

第一項 $\pd{n_x}{x}$ の計算：
$f/g$ の微分は $(f'g - fg')/g^2$ である．
$f=\phi_x$，$g=\sqrt{\phi_x^2+\phi_y^2}$ とすると，
\begin{itemize}
    \item $f' = \phi_{xx}$
    \item $g' = \dfrac{\phi_x\phi_{xx} + \phi_y\phi_{xy}}{\sqrt{\phi_x^2+\phi_y^2}}$
    \quad（連鎖律による）
\end{itemize}
商の微分公式に代入すると：
\begin{equation}
 \pd{n_x}{x}
 = \frac{
     \overbrace{\phi_{xx}}^{f'}
     \overbrace{(\phi_x^2+\phi_y^2)^{1/2}}^{g}
     -
     \overbrace{\phi_x}^{f}
     \overbrace{\left(\dfrac{\phi_x\phi_{xx} + \phi_y\phi_{xy}}{(\phi_x^2+\phi_y^2)^{1/2}}\right)}^{g'}
   }{\underbrace{\phi_x^2+\phi_y^2}_{g^2}}
 \label{eq:dnx_dx_raw}
\end{equation}
通分・整理すると：
\begin{equation}
 \pd{n_x}{x}
 = \frac{\phi_{xx}\phi_y^2 - \phi_x\phi_y\phi_{xy}}{(\phi_x^2+\phi_y^2)^{3/2}}
 \label{eq:dnx_dx}
\end{equation}
同様に $y$ 方向（対称性 $x \leftrightarrow y$）：
\begin{equation}
  \pd{n_y}{y}
  = \frac{\phi_{yy}\phi_x^2 - \phi_y\phi_x\phi_{yx}}{(\phi_x^2+\phi_y^2)^{3/2}}
  \label{eq:dny_dy}
\end{equation}
$\phi_{xy}=\phi_{yx}$（Schwarz の定理）として
式~\eqref{eq:dnx_dx} と式~\eqref{eq:dny_dy} を加算すると，
分子は $\phi_{xx}\phi_y^2 - 2\phi_x\phi_y\phi_{xy} + \phi_{yy}\phi_x^2$ となり，
式~\eqref{eq:curvature_2d} が得られる．

% 円検証は §~\ref{sec:verify_curvature} の数値テストに統合．
